\chapter{Lenguaje y Comunicación}
\section{Comprensión lectora: ámbito de estudio}\label{sec:l1}
\objetivo{Identificar información, interpretar relaciones y evaluar argumentos en reseñas académicas y artículos de investigación.}
\subsection{Leer para responder con evidencia}
En esta sección no se evalúa memorizar definiciones disciplinares. La tarea es comprender el texto proporcionado. Identificar significa localizar datos que satisfacen un criterio; interpretar integra relaciones y sentido; evaluar conecta contenido y forma con una valoración justificada. Las tres operaciones pueden aparecer en un mismo género \cite[págs. 18--23]{ceneval}.

Una estrategia útil consta de lectura global, identificación del propósito y regreso a pasajes concretos según la pregunta. Señale quién afirma cada idea, qué se observó, con qué método y qué conclusión se propone. Una hipótesis no equivale a un resultado; una correlación descrita no demuestra causalidad. Palabras como «podría», «en esta muestra» y «no basta» delimitan alcance.

Una idea central debe abarcar la intención del texto sin incluir afirmaciones nuevas. Un detalle verdadero puede ser distractor si no responde la pregunta global. Una inferencia válida debe apoyarse en señales del texto; no es licencia para completar con opiniones personales. Para evaluar un argumento se revisa si la evidencia es pertinente, si sostiene la conclusión y si se reconocen límites.

\subsection{Reseña académica y artículo de investigación}
La reseña presenta una obra y la valora: puede resumir su aporte, explicar enfoque y señalar límites. Distinga la voz del autor de la obra y la del reseñista. El artículo de investigación organiza un problema, método, resultados y discusión; el esquema concreto varía. En ambos, identifique qué afirmaciones están respaldadas y cuáles son propuestas de trabajo futuro.

\ejemplo{L1}{Reseña: aporte y limitación}
Microtexto original: «El manual ofrece ejemplos claros de estadística aplicada. No explica cómo se eligieron los casos; por eso, su utilidad didáctica no demuestra que cubra toda situación profesional». Aporte: claridad de ejemplos. Límite: selección no explicada. Una síntesis adecuada conserva ambos; «el manual no sirve» elimina el reconocimiento y «resuelve cualquier caso» contradice la cautela. Es un ejemplo del género reseña, no una evaluación de una obra real.
\ejemplo{L2}{Artículo: resultado frente a generalización}
Microtexto original ficticio: «De veinte participantes, quince completaron una tarea con la interfaz A. El estudio se realizó en un solo laboratorio y no comparó grupos de edad». El dato identificable es 15 de 20; la proporción es 75\%. No permite afirmar que la interfaz sea mejor que otra no estudiada ni que todas las edades respondan igual. La evaluación pertinente pide ampliar condiciones antes de generalizar.
\errores{Elegir un dato verdadero que no responde la pregunta; confundir propuesta con hallazgo; atribuir al texto certezas que presenta como posibilidades; ignorar negaciones y límites.}
\sintesis{Propósito, afirmación, evidencia y alcance: toda respuesta debe poder vincularse con uno o más pasajes.}
\fuentes{ceneval,reading}

\section{Comprensión lectora: ámbito literario}\label{sec:l2}
\objetivo{Interpretar acciones, voz, imágenes y estructura en cuentos y ensayos sin exigir una lectura literal de toda expresión.}
\subsection{Narración, perspectiva y sentido}
Un cuento organiza acciones, personajes, tiempo y espacio mediante una voz narrativa. El narrador no debe identificarse automáticamente con el autor. La perspectiva limita qué se conoce y qué se omite; esas omisiones pueden obligar al lector a inferir motivos a partir de acciones y detalles. Distinga el orden en que se cuentan los hechos del orden cronológico.

Tema es una dimensión de sentido, como espera, pérdida o pertenencia; argumento es la secuencia concreta. Un objeto puede adquirir valor simbólico por repetición o posición en el relato. La interpretación más sólida explica varias señales y evita inventar hechos. Un final abierto permite incertidumbre; no convierte todas las conclusiones en igualmente justificadas.

El ensayo literario desarrolla una reflexión con voz personal, asociaciones y recursos expresivos. Puede argumentar sin adoptar el formato de un artículo experimental. Una metáfora compara o traslada sentido; su interpretación debe responder a la función en el pasaje. Preguntar por el efecto de una imagen exige explicar qué aporta, no solo ponerle un nombre técnico \cite[págs. 19--21]{ceneval}.

\ejemplo{L3}{Cuento: inferir desde acciones}
Microcuento original: «Después de vender la casa, Julia guardó una llave que ya no abría ninguna puerta. Cada domingo la cambiaba de bolsillo». El dato literal es conservar y mover la llave. Una interpretación razonable relaciona el objeto con el vínculo con la casa y el hábito de recordar. No se puede concluir que Julia haya recuperado legalmente la vivienda: el relato no lo dice. La diferencia entre inferencia y hecho evita sobreinterpretar.
\ejemplo{L4}{Ensayo: explicar una imagen}
Microensayo original: «El silencio de una biblioteca no es ausencia: es una mesa donde se sientan voces que no coinciden en el tiempo». La imagen propone que la lectura reúne autores y lectores de épocas distintas. No afirma que haya personas físicamente hablando en todas las mesas. Una respuesta evaluativa puede señalar que la paradoja transforma una aparente falta de sonido en presencia simbólica.
\errores{Confundir narrador y autor; convertir símbolos en hechos materiales; imponer una moraleja no sustentada; elegir una interpretación que contradice acciones explícitas.}
\sintesis{Lea tanto lo que ocurre como la forma de contarlo; explique el sentido usando detalles del texto, no datos imaginados.}
\fuentes{ceneval,reading}

\section{Comprensión lectora: participación social}\label{sec:l3}
\objetivo{Aplicar requisitos y procedimientos de convocatorias e interpretar notas informativas distinguiendo hechos, fuentes y decisiones pendientes.}
\subsection{Convocatorias y notas}
Una convocatoria comunica quién invita, a quién se dirige, qué se ofrece, requisitos, plazos, procedimiento y criterios de selección. Para responder construya una lista de condiciones simultáneas. Cumplir edad pero incumplir extensión o fecha no basta. No invente requisitos por parecer habituales: si el texto no exige experiencia, no la convierta en criterio.

Una nota informativa distingue hechos reportados, atribuciones, antecedentes y proyecciones. Identifique quién dijo qué, fecha, lugar, alcance y estado de una decisión. «Se estudiará ampliar» no equivale a «se amplió». Una valoración de una persona debe conservar su atribución; no se presenta como consenso de toda la población.

En evaluación, pregunte qué información falta para valorar una decisión o qué función cumple un dato. La respuesta debe ser pertinente al propósito del texto, sin exigir detalles irrelevantes. Los ejemplos del banco son ficticios y no anuncian trámites o eventos reales.
\ejemplo{L5}{Convocatoria: condiciones acumulativas}
Microconvocatoria original: «Personas mayores de 18 años podrán solicitar uno de diez lugares. Enviarán un texto propio de 200 a 300 palabras antes del viernes. Si hay más solicitudes válidas, habrá sorteo». Una persona de 20 años con 250 palabras a tiempo cumple; una de 17 no cumple edad; una de 25 con 500 palabras incumple extensión. El sorteo se activa solo si se supera el cupo con solicitudes válidas.
\ejemplo{L6}{Nota: alcance de una decisión}
Micronota original: «La dirección anunció una prueba de horario extendido durante dos semanas. Después medirá asistencia y decidirá si continúa». El horario se aplicará como prueba, no como política permanente ya decidida. Una pregunta sobre continuidad se responde con el estado pendiente. Para evaluar el resultado sería pertinente conocer asistencia y criterios, no suponer éxito por el anuncio.
\errores{Confundir publicación de resultados con fecha de inscripción; omitir una condición; tomar una cita como afirmación del redactor; interpretar un piloto como programa definitivo.}
\sintesis{Convierta requisitos en una lista; separe hecho, declaración, condición y decisión futura.}
\fuentes{ceneval,reading}

\section{Redacción indirecta: ámbito de estudio}\label{sec:e1}
\objetivo{Elegir textos académicos adecuados al propósito, coherentes y conformes con la norma del español.}
\subsection{Dimensión comunicativa}
La prueba pide seleccionar la formulación adecuada, no escribir un ensayo completo. La tabla operativa de la página 30 del PDF enumera artículo de divulgación científica, protocolo de proyecto de investigación y reseña en estudio; editorial de periódico, convocatoria y carta de exposición de motivos en participación social. Los descriptores de la página 22 también mencionan ensayo académico y noticia. Se conserva esa distinción entre tabla y descriptores, sin sustituir una por otra. El propósito determina qué información debe incluirse: un objetivo indica qué se hará; el método cómo; los resultados qué se observó; la discusión qué significa y con qué límites. Un registro académico suele favorecer precisión, atribución de fuentes y lenguaje explícito, sin convertir toda primera persona en error automático.

Una conclusión debe corresponder a la evidencia. «Se observó asociación en esta muestra» no significa «se demostró causalidad universal». Evite cambiar certeza, población o temporalidad al reformular. La fidelidad semántica tiene prioridad sobre que una frase suene más enfática \cite[págs. 29--31]{ceneval}.

\subsection{Dimensión gramatical y semántica}
La concordancia conecta sujeto y verbo y sustantivos con modificadores. Un sujeto coordinado suele requerir plural; construcciones especiales deben interpretarse según su estructura. Un referente como «su propuesta» puede resultar ambiguo si hay dos posibles titulares. Repetir un nombre puede ser mejor que una sustitución elegante pero confusa.

Los conectores expresan causa, consecuencia, contraste, adición o explicación. «Por ello» requiere un antecedente causal; «sin embargo» introduce contraste; «es decir» reformula. No se eligen solo por sonar formales. La cohesión relaciona oraciones; la coherencia conserva sentido y progresión global.

El haber existencial es impersonal: «hubo problemas», «había tres solicitudes». No debe confundirse con el auxiliar en «habían terminado», donde la estructura es otra. Evite dequeísmo cuando un verbo introduce directamente «que», como «confirmó que»; no suprima la preposición cuando sí la exige, como en «se dio cuenta de que» \cite[entradas concordancia, haber, queísmo y dequeísmo]{rae}.

\subsection{Norma de referencia y dimensión ortográfica}
La página 31 de la guía oficial indica que parte de la prueba se basa en el DPD de 2005 y excluye segundas ediciones o actualizaciones para esa evaluación. Esta guía toma esa edición como referencia de examen y utiliza casos básicos estables entre versiones. La consulta en línea actual se identifica separadamente y no se usa para imponer novedades a los reactivos \cite{raehoy}. Ante una diferencia, prevalece para la preparación del examen la referencia expresamente señalada en el PDF aportado \cite[pág. 31]{ceneval}.

La tilde distingue palabras por reglas generales y casos diacríticos: tú pronombre y tu posesivo; él pronombre y el artículo; sé verbal y se pronombre. Los interrogativos y exclamativos llevan tilde también en usos indirectos: «explicó cómo se midió». Aún equivale a todavía; aun suele equivaler a incluso. El contexto, no la mera posición, decide.

La coma separa elementos e incisos y delimita vocativos; no se pone entre sujeto y predicado por la sola longitud del sujeto. Dos puntos pueden anunciar una enumeración cuando la construcción está completa. La puntuación debe preservar relaciones sintácticas y sentido. Una coma no reemplaza cualquier pausa de respiración \cite[entradas tilde, coma, dos puntos, aun]{rae}.
\ejemplo{E1}{Adecuación y precisión}
Para un protocolo: «Comparar el tiempo de respuesta de dos interfaces en la tarea de registro». Es mejor que «mostrar que nuestra interfaz es maravillosa»: define una operación y variable sin anticipar el resultado. Para verificarlo faltaría delimitar participantes, condiciones y métrica según el estudio.
\ejemplo{E2}{Concordancia y referencia}
«La hipótesis y el procedimiento fueron revisados» concuerda con un sujeto plural de géneros distintos. En «Ana informó a Lucía que su propuesta fue aceptada», su no identifica inequívocamente a la titular. «La propuesta de Ana» resuelve el referente si esa es la intención. Corregir no significa alterar de quién era la propuesta.
\ejemplo{E3}{Acentuación y puntuación}
«El informe explicó cómo se obtuvo la muestra; aún faltan dos anexos». El es artículo sin tilde; cómo es interrogativo indirecto; aún equivale a todavía. El punto y coma une segmentos relacionados pero con autonomía sintáctica. En cambio, «Los resultados de la prueba, muestran mejoras» separa indebidamente sujeto y verbo.
\errores{Seleccionar una opción por su tono solemne aunque cambie el sentido; quitar todas las preposiciones antes de que; usar coma entre sujeto y verbo; extrapolar resultados limitados.}
\sintesis{Verifique propósito, significado, sintaxis y ortografía en ese orden, volviendo al contexto para resolver ambigüedades.}
\fuentes{ceneval,rae,writing}

\section{Redacción indirecta: participación social}\label{sec:e2}
\objetivo{Seleccionar avisos, noticias, convocatorias y cartas de motivos claros, pertinentes y correctamente escritos.}
\subsection{Propósito, destinatario y registro}
Un editorial de periódico sostiene una postura institucional mediante argumentos; se distingue de la nota informativa, cuyo propósito central es informar. Un artículo de divulgación explica conocimientos para lectores no especialistas, con precisión y ejemplos accesibles. Una reseña presenta y valora una obra con criterios.

Una solicitud debe decir qué se pide, a quién y con qué fundamento. Una carta de motivos relaciona antecedentes y objetivos con la oportunidad solicitada; no presupone una resolución que aún no existe. Un aviso prioriza acción, lugar y plazo. Una noticia informa y atribuye valoraciones, en lugar de presentarlas como hechos indiscutibles.

El registro depende de relación y contexto. Una petición institucional suele requerir respeto y precisión; no basta añadir fórmulas ceremoniosas si falta el trámite concreto. La claridad beneficia al destinatario: «entregue el formulario firmado en recepción antes de las 16:00 del viernes» comunica más que «haga eso pronto». Si cambia solo la sede, indíquelo y confirme las condiciones que siguen vigentes para evitar inferencias.

\subsection{Gramática y ortografía en contexto}
Conserve concordancia, referencias claras y paralelismo en enumeraciones: «presentar identificación, completar el formulario y firmar la solicitud». Para negar una alternativa y contraponer otra se usa sino: «no se canceló, sino que cambió de sede»; si no expresa condición negativa: «si no asiste, avise». Porque introduce causa; por qué aparece en preguntas directas o indirectas; el porqué es sustantivo; por que depende de construcciones específicas.

Ha es auxiliar en «ha enviado»; a es preposición en «a recepción». Hay expresa existencia, ahí lugar y ay una interjección. El vocativo se separa con coma: «María, revise el aviso». La selección correcta requiere leer la oración completa, no aplicar una regla por una palabra aislada \cite[entradas sino, porque, haber, ahí y ay]{rae}.
\ejemplo{E1}{Solicitud pertinente}
«Solicito la reconsideración de mi expediente porque se incorporó la constancia pendiente. Adjunto el documento y agradezco la revisión». La petición, razón y evidencia están relacionadas. «Confirmen que gané» no corresponde a una revisión pendiente; «ya saben lo que quiero» omite información necesaria.
\ejemplo{E2}{Nota con atribución}
«La portavoz calificó la medida de insuficiente» informa una valoración y su fuente. «La medida es inútil y todos lo saben» convierte un juicio en consenso no demostrado. Para mantener fidelidad, no borre quién habla ni cambie una opinión por una decisión institucional.
\ejemplo{E3}{Condición y contraste}
«No se exige experiencia, sino disponibilidad; si no puede asistir, avise». El primer sino contrapone requisitos; el segundo si no introduce una condición negativa. «Aún hay lugares» significa que todavía quedan; «aun sin experiencia» equivale a incluso sin ella. La tilde responde al significado.
\errores{Elegir un mensaje formal pero impreciso; atribuir decisiones pendientes como definitivas; confundir condición y contraste; omitir plazos o destinatarios; corregir ortografía alterando el contenido.}
\sintesis{Una redacción adecuada comunica exactamente lo necesario a su destinatario, con relaciones gramaticales y signos que preservan el sentido.}
\fuentes{ceneval,rae,writing}
